\documentclass[tikz,border=6pt]{standalone}
\usepackage{ctex}
\usepackage{amsmath}
\usetikzlibrary{calc, arrows.meta}

\begin{document}
\begin{tikzpicture}[scale=1.1, thick]

% 通径示意图：过焦点垂直对称轴的弦，长度 = 2p
% 抛物线 y^2 = 4x (p=2)，焦点 F(1,0)，通径两端 C(1,2), D(1,-2)

% 坐标轴
\draw[->, thin, gray] (-2.0, 0) -- (5.0, 0) node[right, font=\small] {$x$};
\draw[->, thin, gray] (0, -3.2) -- (0, 3.2) node[above, font=\small] {$y$};
\node[font=\footnotesize, below right] at (0,0) {$O$};

% 抛物线
\draw[black, thick, domain=-2.8:2.8, smooth, samples=80]
  plot ({(\x)^2/4}, \x);

% 准线 x = -1
\draw[red, thin] (-1, -3.0) -- (-1, 3.0);
\node[red, font=\small, above] at (-1, 3.0) {$x=-\dfrac{p}{2}$};

% 焦点 F
\fill[black] (1, 0) circle [radius=2.5pt];
\node[below right, font=\small] at (1, 0) {$F\!\left(\dfrac{p}{2},0\right)$};

% 通径端点 C(1,2), D(1,-2)（将 x=1 代入 y^2=4x => y=±2）
\coordinate (C) at (1, 2);
\coordinate (D) at (1, -2);

% 通径线段（红色实线，突出）
\draw[red, very thick] (C) -- (D);

\fill[red] (C) circle [radius=2.5pt];
\fill[red] (D) circle [radius=2.5pt];

\node[right, font=\small] at (C) {$C\!\left(\dfrac{p}{2}, p\right)$};
\node[right, font=\small] at (D) {$D\!\left(\dfrac{p}{2}, -p\right)$};

% 双箭头标注通径长度
\draw[red, <->, thick] (1.6, -2) -- (1.6, 2)
  node[midway, right, font=\small] {通径 $= 2p$};

% 垂直于轴的标记（通径 ⊥ 对称轴）
\draw[thin] (1, 0.18) -- (1-0.18, 0.18) -- (1-0.18, 0);

% C, D 到准线的射影（灰色虚线）
\draw[gray, dashed, thin] (C) -- (-1, 2);
\draw[gray, dashed, thin] (D) -- (-1, -2);
\fill[gray] (-1, 2) circle [radius=1.5pt];
\fill[gray] (-1, -2) circle [radius=1.5pt];

% 标注 p/2 距离
\draw[gray, <->, thin] (-1, 2.5) -- (1, 2.5)
  node[midway, above, font=\footnotesize, gray] {$p$};

% 公式文本框
\node[font=\small, align=center,
      draw=red!60, rounded corners=3pt, fill=red!5,
      inner sep=5pt] at (3.5, -2.4)
  {通径 $= 2p$\\（过焦点 $\perp$ 对称轴的弦）};

\end{tikzpicture}
\end{document}
